\documentclass{inf-ufg}
\usepackage[utf8]{inputenc}
\usepackage[T1]{fontenc}
\usepackage{geometry}
\usepackage{graphicx}
\usepackage{float}
\usepackage{hyperref}
\usepackage[english,indonesian]{babel}

\title{Optimasi Pemrosesan Template LaTeX Menggunakan Pendekatan Berbasis AI}
\author{Academic Support Agent}
\date{\today}

\begin{document}
\maketitle

\begin{abstract}

\end{abstract}

\chapter{PENDAHULUAN}


\chapter{TINJAUAN PUSTAKA}


% Original content (commented out):
% % Exemplo de dissertação do INF-UFG com texto em portugues formatado com LaTeX
\documentclass[tese, abnt]{inf-ufg}
% Opções da classe inf-ufg (ao usar mais de uma, separe por vírgulas)
%   [tese]         -> Tese de doutorado.
%   [dissertacao]  -> Dissertação de mestrado (padrão).
%   [monografia]   -> Monografia de especialização.
%   [relatorio]    -> Relatório final de graduação.
%   [abnt]         -> Usa o estilo "abnt-alf" de citação bibliográfica.
%   [nocolorlinks] -> Os links de navegação no texto ficam na cor preta.
%                     Use esta opção para gerar o arquivo para impressão
%                     da versão final do seu texto!!!

%----------------------------------------------------- INICIO DO DOCUMENTO %
\begin{document}

%------------------------------------------ AUTOR, TÍTULO E DATA DE DEFESA %
\autor{Jose da Silva}
\autorR{Silva, Jose}

\titulo{\textless Título do Trabalho\textgreater}
\subtitulo{\textless Subtítulo do Trabalho\textgreater}

\cidade{Goiânia}
\dia{02}
\mes{03}
\ano{2017} % Formato numérico: \dia{01}, \mes{01} e \ano{2009}

%-------------------------------------------------------------- ORIENTADOR %
\orientador{Abrãao Santana}
\orientadorR{Santana, Abrãao}
% Use os comandos a seguir se for Orientadora e nao Orientador.
%\orientadora{\textless Nome da Orientadora\textgreater}
%\orientadoraR{\textless Nome Reverso da Orientadora\textgreater}

%\coorientador{\textless Nome do Co-orientador\textgreater}
%\coorientadorR{\textless Nome Reverso do Co-orientador\textgreater}
% Use os comandos a seguir se for Co-orientadora e nao Coorientador.
%\coorientadora{\textless Nome da Co-orientadora\textgreater}
%\coorientadoraR{\textless Nome Reverso da Co-orientadora\textgreater}

%-------------------------------------------------- INSTITUIÇÃO E PROGRAMA %
\universidade{Universidade Federal de Goiás}
\uni{UFG}
\unidade{Insituto de Informática}
\departamento{} % Unidades com mais de um departamento.

%\universidadeco{\textless Nome da Universidade do Co-orientador\textgreater}
%\unico{\textless Sigla da Universidade do Co-orientador\textgreater}
%\unidadeco{\textless Nome da Unidade Acadêmica do Co-orientador\textgreater}

\programa{Ciência da Computação}
\concentracao{Ciência da Computação}

%-------------------------------------------------- ELEMENTOS PRÉ-TEXTUAIS %
\capa    % Gera o modelo da capa externa do trabalho
\publica % Gera a autorização para publicação em formato eletrônico
\rosto   % Primeira folha interna do trabalho

%\input{./pre/pre_biblioteca}  %Ficha Catalográfica expedida pela secretaria do programa
\input{./pre/pre_aprovacao}
\input{./pre/pre_direitos}
\input{./pre/pre_dedicatoria}
\input{./pre/pre_agradecimentos}
\input{./pre/pre_epigrafe}
\input{./pre/pre_resumo}
\input{./pre/pre_abstract}

\tabelas[figtabalgcod]
%Opções:
%nada [] -> Gera apenas o sumário
%fig     -> Gera o sumário e a lista de figuras
%tab     -> Sumário e lista de tabelas
%alg     -> Sumário e lista de algoritmos
%cod     -> Sumário e lista de códigos de programas
%
% Pode-se usar qualquer combinação dessas opções.
% Por exemplo:
%  figtab       -> Sumário e listas de figuras e tabelas
%  figtabcod    -> Sumário e listas de figuras, tabelas e
%                  códigos de programas
%  figtabalg    -> Sumário e listas de figuras, tabelas e algoritmos
%  figtabalgcod -> Sumário e listas de figuras, tabelas, algoritmos e
%                  códigos de programas

%--------------------------------------------------------------- CAPÍTULOS %
\input{./tex/cap_I}
%\input{./tex/cap_II}
%\input{./tex/cap_III}

%------------------------------------------------------------ BIBLIOGRAFIA %
\cleardoublepage
%\nocite{*} %%% Retire esta linha para gerar a bibliografia com apenas as
           %%% referências usadas no seu texto!
\arial
\bibliography{./bib/tese} %%% Nomes dos seus arquivos .bib
\label{ref-bib}

%--------------------------------------------------------------- APÊNDICES %
\apendices

\input{./pos/apend_I}
%\input{./pos/apend_II}

\end{document}

%------------------------------------------------------------------------- %
%        F I M   D O  A R Q U I V O :  m o d e l o - t e s e . t e x       %
%------------------------------------------------------------------------- %


\end{document}